\documentclass{report}
\usepackage[utf8]{inputenc}
\usepackage[T1]{fontenc}
\usepackage[french]{babel}
\usepackage{geometry}
\usepackage{fancyhdr}
\usepackage{lastpage}
\usepackage{titlesec}
\usepackage{listings}
\usepackage{enumitem}
\usepackage{color}
\usepackage{xcolor}
\usepackage{soul}
\usepackage{graphicx}
\usepackage[utf8]{inputenc}
\usepackage{amsmath}

\newcommand{\DocName}{Projet Microcontrôleurs - Rapport}
\newcommand{\DocAuthors}{\bsc{Da Silva} Tristan - \bsc{Battaglia} Charles-Antoine \\ \bsc{Botté-Magalhaes} Jules - \bsc{Gauthier} Keanu}

\definecolor{dkgreen}{rgb}{0,0.6,0}
\definecolor{gray}{rgb}{0.5,0.5,0.5}
\definecolor{mauve}{rgb}{0.58,0,0.82}
\definecolor{addrred}{rgb}{0.93,0.43,0.43}
\definecolor{commentgray}{rgb}{0.45,0.45,0.45}

\lstdefinelanguage{MicroAsm}{
    sensitive=true,
    morecomment=[l]{\#},
    moredelim=[s][\color{addrred}]{@}{:}
}

\lstdefinelanguage{HexProg}{
    sensitive=true,
    morecomment=[l]{;}
}

\lstdefinelanguage{AsmStd}{
    sensitive=true,
    morecomment=[l]{;},
    morecomment=[l]{\#},
    morekeywords={
        MOV,ADD,SUB,MUL,DIV,INC,DEC,CMP,AND,OR,XOR,NOT,
        JMP,JE,JNE,JG,JL,JGE,JLE,CALL,RET,PUSH,POP,
        LDA,STA,LOAD,STORE,INPUT,OUTPUT,HALT,NOP
    }
}

\lstset{frame=tb,
    language=MicroAsm,
    aboveskip=3mm,
    belowskip=3mm,
    showstringspaces=false,
    columns=flexible,
    basicstyle={\small\ttfamily},
    numbers=none,
    numberstyle=\tiny\color{gray},
    keywordstyle=\color{blue},
    commentstyle=\color{commentgray}, 
    breaklines=true,
    breakatwhitespace=true,
    tabsize=3,
   }

\lstset{breaklines=true}

\geometry{a4paper, margin=1in}

% Suppression de l'affichage du mot "Chapitre"
\titleformat{\chapter}[display]
    {\normalfont\bfseries\LARGE}
    {} % Supprime le label "Chapitre X"
    {0pt} % Espacement entre le label et le titre
    {}

\definecolor{Light}{gray}{.90}
\sethlcolor{Light}

\let\OldTexttt\texttt
\renewcommand{\texttt}[1]{\OldTexttt{\hl{#1}}}

\renewcommand{\thechapter}{\arabic{chapter}}
\renewcommand{\thesection}{\thechapter.\arabic{section}}
\renewcommand{\thesubsection}{\thesection.\arabic{subsection}}

\fancyhf{}
\fancyhead[L]{\DocName}
\fancyfoot[L]{\DocAuthors}
\fancyfoot[R]{\thepage / \pageref{LastPage}}
\renewcommand{\headrulewidth}{0 pt}
\renewcommand{\footrulewidth}{0.2pt}
\pagestyle{fancy}
\fancypagestyle{plain}{
    \fancyhf{}
    \fancyfoot[R]{\thepage / \pageref{LastPage}}
    \fancyfoot[L]{\DocAuthors}
    \fancyhead[L]{\DocName}
}

\title{\DocName}
\author{\DocAuthors}
\date{\textbf{JUNIA - ISEN - AP3}}

\newcounter{questioncounter}
\setcounter{questioncounter}{0}
\begin{document}
\maketitle
\tableofcontents

\chapter{Calculs valeurs analogiques}

\section*{Dimensionnement du FILTER1}
On a :
\[f = 250 \text{ Hz}
\]
\[\omega = 2\pi f = 2\pi \times 250 = 500\pi \text{ rad/s}
\]

La fonction de transfert \( H_1(j\omega) \) est donnée par :
\[H_1(j\omega) = \frac{1}{1 + j R_0 C_0 \omega}
\]

On définit la fréquence de coupure \( \omega_c \) par :
\[
\omega_c = \frac{1}{R_0 C_0}
\]
\[
f_c = \frac{\omega_c}{2\pi} = \frac{1}{2\pi R_0 C_0}=250 \text{ Hz}
\]

On fixe :
\[C_0 = 47 \times 10^{-9} \text{ F}
\]
\[R_0 = \frac{1}{2\pi \times 250 \times 47 \times 10^{-9}} = 13,5 \text{ k}\Omega
\]

On prend donc pour valeurs de composants à disposition :
\[
\fbox{$R_0 = 10 \text{ k}\Omega$}
\]
\[
\fbox{$C_0 = 47 \times 10^{-9} \text{ F}$}
\]


\newpage
\section*{Dimensionnement du FILTER2}

On a :
\[
f = 4 \times 10^3 \text{ Hz}
\]
\[
\omega = 2\pi f = 2\pi \times 4 \times 10^3 = 8000\pi \text{ rad/s}
\]

La fonction de transfert \( H_2(j\omega) \) est donnée par :
\[
H_2(j\omega) = \frac{j R_1 C_1 \omega}{1 + j R_1 C_1 \omega}
\]

On définit la fréquence de coupure \( \omega_c \) par :
\[
\omega_c = \frac{1}{R_1 C_1} \quad \Leftrightarrow \quad \frac{\omega}{\omega_c} = R_1 C_1 \omega
\]

Donc :
\[
H_2(j\omega) = \frac{1}{1 + j \frac{\omega}{\omega_c}} \times j \frac{\omega}{\omega_c} = \frac{j \frac{\omega}{\omega_c}}{1 + j \frac{\omega}{\omega_c}}
\]

On fixe :
\[
C_1 = 6,8 \times 10^{-9} \text{ F}
\]
\[
R_1 = \frac{1}{4 \times 10^3 \times 2\pi \times 6,8 \times 10^{-9}} = 5,6 \text{ k} \Omega
\]

On prend donc pour valeurs de composants à disposition :
\[
\fbox{$R_1 = 5,6 \text{ k}\Omega$}
\]
\[
\fbox{$C_1 = 6,8 \times 10^{-9} \text{ F}$}
\]

\newpage
\section*{Dimensionnement du FILTER3}
On a :
\[f_{c3L} = 250 \text{ Hz} \quad \text{et} \quad f_{c3H} = 1 \times 10^3 \text{ Hz}
\] 

La fonction de transfert \( H_3(j\omega) \) est donnée par :
\[H_3(j\omega) = -\frac{R_3}{R_2} \times \frac{j R_2 C_2 \omega}{1 + j R_2 C_2 \omega} \times \frac{1}{1 + j R_3 C_3 \omega}
\]

On pose: 
\[R_3 = R_2 = R
\]

Donc :
\[H_3(j\omega) = -\frac{1}{1 + j R C_2 \omega} \times j R C_2 \omega \times \frac{1}{1 + j R C_3 \omega}
\]

On a : 
\[\omega_{c2} = \frac{1}{R C_2} \quad \text{et} \quad \omega_{c3} = \frac{1}{R C_3}\]

Donc :
\[H_3(j\omega) = -\frac{1}{1 + j \frac{\omega}{\omega_{c2}}} \times j \frac{\omega}{\omega_{c2}} \times \frac{1}{1 + j \frac{\omega}{\omega_{c3}}}
\]

On fixe :
\[R = 10 \text{ k}\Omega
\]

On pose: 
\[\frac{\omega_{c2}}{2\pi} = \frac{1}{2\pi R C_2} = f_{c3L} = 250 \text{ Hz}
 \quad \text{et} \quad \frac{\omega_{c3}}{2\pi} = \frac{1}{2\pi R C_3} = f_{c3H} = 1 \times 10^3 \text{ Hz}
\]

On a donc :
\[C_2 = \frac{1}{2\pi R f_{c3L}} = \frac{1}{2\pi \times 10 \times 10^3 \times 250} = 64 \times 10^{-9} \text{ F}
\]
\[C_3 = \frac{1}{2\pi R f_{c3H}} = \frac{1}{2\pi \times 10 \times 10^3 \times 1 \times 10^3} = 16 \times 10^{-9} \text{ F}
\]

On prend donc pour valeurs de composants à disposition :
\[\fbox{$R_2 = R_3 = 10 \text{ k}\Omega$}
\]
\[\fbox{$C_2 = 68 \times 10^{-9} \text{ F}$}
\]
\[\fbox{$C_3 = 15 \times 10^{-9} \text{ F}$}
\]

\newpage
\section*{Dimensionnement du FILTER4}
On a :
\[f_{c4L} = 1 \times 10^3 \text{ Hz} \quad \text{et} \quad f_{c4H} = 4 \times 10^3 \text{ Hz}
\] 

La fonction de transfert \( H_4(j\omega) \) est donnée par :
\[H_4(j\omega) = -\frac{R_5}{R_4} \times \frac{j R_4 C_4 \omega}{1 + j R_4 C_4 \omega} \times \frac{1}{1 + j R_5 C_5 \omega}
\]

On pose: 
\[R_4 = R_5 = R
\]

Donc :
\[H_4(j\omega) = -\frac{1}{1 + j R C_4 \omega} \times j R C_4 \omega \times \frac{1}{1 + j R C_5 \omega}
\]

On a : 
\[\omega_{c4} = \frac{1}{R C_4} \quad \text{et} \quad \omega_{c5} = \frac{1}{R C_5}\]

Donc :
\[H_4(j\omega) = -\frac{1}{1 + j \frac{\omega}{\omega_{c4}}} \times j \frac{\omega}{\omega_{c4}} \times \frac{1}{1 + j \frac{\omega}{\omega_{c5}}}
\]

On fixe :
\[R = 10 \text{ k}\Omega
\]

On pose: 
\[\frac{\omega_{c4}}{2\pi} = \frac{1}{2\pi R C_4} = f_{c4L} = 1 \times 10^3 \text{ Hz}
 \quad \text{et} \quad \frac{\omega_{c5}}{2\pi} = \frac{1}{2\pi R C_5} = f_{c4H} = 4 \times 10^3 \text{ Hz}
\]

On a donc :
\[C_4 = \frac{1}{2\pi R f_{c4L}} = \frac{1}{2\pi \times 10 \times 10^3 \times 1 \times 10^3} = 16 \times 10^{-9} \text{ F}
\]
\[C_5 = \frac{1}{2\pi R f_{c4H}} = \frac{1}{2\pi \times 10 \times 10^3 \times 4 \times 10^3} = 3,9 \times 10^{-9} \text{ F}
\]

On prend donc pour valeurs de composants à disposition :
\[\fbox{$R_4 = R_5 = 10 \text{ k}\Omega$}
\]
\[\fbox{$C_4 = 15 \times 10^{-9} \text{ F}$}
\]
\[\fbox{$C_5 = 4,7 \times 10^{-9} \text{ F}$}
\]

\end{document}
